\documentclass[11pt]{article}

\usepackage[margin=1in]{geometry}
\usepackage{hyperref}
\usepackage{booktabs}
\usepackage{graphicx}
\usepackage{amsmath,amssymb}
\usepackage{microtype}
\usepackage{xcolor}
\usepackage{natbib}

\title{\textbf{Pareto-Bench}: Cost-Matched Evaluation of\\
Multi-Agent Debate Strategies for LLM Agents}

\author{
  First Author\thanks{Corresponding author: \texttt{email@institution.edu}} \\
  Institution \\
  \and
  Second Author \\
  Institution
}

\date{\today}

\begin{document}

\maketitle

% ============================================================
\begin{abstract}
% ============================================================

Multi-agent debate architectures have been reported to outperform single-agent
baselines on a range of reasoning benchmarks.  However, these comparisons
rarely control for inference cost: a debate pipeline may consume 6--10$\times$ the
tokens of a single call, and it is unclear whether equivalent spend on a
stronger single-agent strategy (e.g., self-consistency with additional samples)
would close the gap.  We introduce \textbf{Pareto-Bench}, a benchmark that
enforces an \emph{iso-cost protocol}: every strategy variant is evaluated under
identical USD token budgets, enabling fair comparison via Pareto frontiers over
(accuracy, cost) space.  We evaluate five strategies---single-agent
self-consistency, vanilla debate~\citep{du2024improving}, role-based debate,
heterogeneous-model debate, and a reasoning-bank baseline~\citep{ouyang2025reasoningbank}---across four benchmarks (GSM8K, ARC-Challenge, GAIA Level 1,
HumanEval) and three budget tiers.  Our results suggest that \ldots
[TODO: fill in after running experiments].

\end{abstract}

% ============================================================
\section{Introduction}
% ============================================================

Large language model (LLM) agents have shown strong performance on complex
reasoning tasks, particularly when multiple agents collaborate through debate
or critique cycles~\citep{du2024improving,cemri2025mast}.  Yet a fundamental
confound persists in the literature: debate pipelines consume substantially more
tokens than single-call baselines.  Whether the gains are attributable to the
\emph{architecture} or merely the \emph{additional computation} is largely
unexplored.

Self-consistency~\citep{wang2023selfconsistency} provides a natural budget-matched
baseline: rather than running three agents for two rounds, one can draw an
equivalent number of independent samples from a single model and take the
majority vote.  On many tasks, this comparison has not been made explicitly.

We address this gap with \textbf{Pareto-Bench}, a reusable benchmark framework
that:
\begin{enumerate}
    \item enforces a shared USD cost cap across all strategy variants;
    \item evaluates strategies on four diverse benchmarks spanning math
          reasoning, multiple-choice science, open-ended agentic tasks, and
          code generation;
    \item draws Pareto frontiers over (accuracy, cost) space rather than
          reporting flat accuracy tables.
\end{enumerate}

% ============================================================
\section{Related Work}
% ============================================================

\paragraph{Multi-agent debate.}
\citet{du2024improving} propose a homogeneous debate procedure in which agents
iteratively refine answers after observing peers' responses.  \citet{cemri2025mast}
provide a taxonomy of multi-agent systems and identify debate as a prominent
coordination pattern.

\paragraph{Self-consistency.}
\citet{wang2023selfconsistency} show that sampling multiple chain-of-thought
reasoning traces and taking the majority vote substantially improves accuracy
on arithmetic and commonsense benchmarks.  This serves as the primary
iso-cost baseline in our study.

\paragraph{Memory-augmented agents.}
\citet{ouyang2025reasoningbank} augment single-agent inference with a bank of
stored reasoning traces retrieved at inference time.  We include this as a
fifth strategy to test whether memory access can close the gap with debate at
matched cost.

\paragraph{Benchmark evaluation.}
Existing multi-agent benchmarks do not control for inference cost.
Pareto-Bench is the first systematic effort to compare strategies on an
iso-cost basis across multiple task domains.

% ============================================================
\section{Method: The Iso-Cost Protocol}
% ============================================================

\subsection{Cost accounting}

All API calls pass through a \texttt{CostTracker} that records
$(m, n_{\text{in}}, n_{\text{out}}, c, t, s, \tau, r)$---model, input tokens,
output tokens, cost in USD, timestamp, strategy, task ID, and run ID---to a
JSONL log.  Costs are computed from a versioned model registry (see
\texttt{configs/models.yaml}).

\subsection{Budget caps}

Each experiment cell $(s, d, B, \xi)$ for strategy $s$, dataset $d$, budget
$B$, and seed $\xi$ receives a hard USD cap enforced by
\texttt{BudgetExceededError}.  When the cap is hit, the run terminates cleanly
and partial results are recorded.  This ensures that every strategy completes
approximately the same number of tasks under the same total spend.

\subsection{Pareto evaluation}

Let $\mathcal{S} = \{(c_s, a_s)\}$ be the set of (mean cost, mean accuracy)
points over strategies.  A strategy $s$ dominates $s'$ if
$c_s \le c_{s'}$ and $a_s \ge a_{s'}$ with at least one strict inequality.
The Pareto frontier is the set of non-dominated strategies.

% ============================================================
\section{Experimental Setup}
% ============================================================

\subsection{Strategies}

\begin{itemize}
    \item \textbf{SingleAgentSC}: $N$ independent samples, majority vote
          (self-consistency baseline; \citealt{wang2023selfconsistency}).
    \item \textbf{VanillaDebate}: $A$ agents, $R$ rounds, majority vote
          (\citealt{du2024improving}).
    \item \textbf{RoleDebate}: Solver / Critic / Judge pipeline, $R$ cycles.
    \item \textbf{HeteroDebate}: Mixed model families in each role.
    \item \textbf{ReasoningBank}: Single agent with in-context retrieval
          from accumulated traces (\citealt{ouyang2025reasoningbank}).
\end{itemize}

\subsection{Benchmarks}

Table~\ref{tab:benchmarks} summarises the four benchmarks.

\begin{table}[h]
\centering
\caption{Benchmark summary.}
\label{tab:benchmarks}
\begin{tabular}{llll}
\toprule
\textbf{Benchmark} & \textbf{Task type} & \textbf{Metric} & \textbf{\# test tasks} \\
\midrule
GSM8K              & Math word problems & Exact-match (number) & 1,319 \\
ARC-Challenge      & Science MCQ        & Letter accuracy       & 1,172 \\
GAIA Level 1       & Agentic reasoning  & Exact-match           & TBD \\
HumanEval          & Code generation    & pass@1 (functional)   & 164 \\
\bottomrule
\end{tabular}
\end{table}

\subsection{Models}

All experiments use \texttt{gpt-4o-mini} as the default model.
Heterogeneous debate uses \texttt{gpt-4o-mini} (solver),
\texttt{claude-haiku-4-5} (critic), and \texttt{gemini-2.5-flash} (judge).

\subsection{Budget tiers}

We evaluate at three tiers: \$0.50, \$2.00, and \$5.00 per run.
Three seeds (42, 123, 7) are used for variance estimation.

% ============================================================
\section{Results}
% ============================================================

[TODO: fill in after running experiments.  Reference Figure~\ref{fig:headline},
Figure~\ref{fig:hetero}, Figure~\ref{fig:per_dataset}, and Table~\ref{tab:main}.]

\begin{figure}[h]
\centering
\includegraphics[width=0.55\textwidth]{figures/fig1_headline_pareto.pdf}
\caption{Headline Pareto frontier: mean success rate vs.\ mean cost per task,
aggregated across all four benchmarks.}
\label{fig:headline}
\end{figure}

\begin{figure}[h]
\centering
\includegraphics[width=0.7\textwidth]{figures/fig2_heterogeneity.pdf}
\caption{Heterogeneity effect: vanilla debate (homogeneous) vs.\
heterogeneous debate (mixed model families) per dataset.}
\label{fig:hetero}
\end{figure}

\begin{figure}[h]
\centering
\includegraphics[width=\textwidth]{figures/fig3_per_dataset.pdf}
\caption{Per-dataset breakdown of mean success rate across strategies.}
\label{fig:per_dataset}
\end{figure}

\begin{table}[h]
\centering
\caption{Main results: mean success rate ± std (strategy × dataset × \$2 budget tier).}
\label{tab:main}
\begin{tabular}{lcccc}
\toprule
\textbf{Strategy} & \textbf{GSM8K} & \textbf{ARC} & \textbf{GAIA-L1} & \textbf{HumanEval} \\
\midrule
SingleAgentSC  & -- & -- & -- & -- \\
VanillaDebate  & -- & -- & -- & -- \\
RoleDebate     & -- & -- & -- & -- \\
HeteroDebate   & -- & -- & -- & -- \\
ReasoningBank  & -- & -- & -- & -- \\
\bottomrule
\end{tabular}
\end{table}

% ============================================================
\section{Discussion}
% ============================================================

[TODO]

% ============================================================
\section{Limitations}
% ============================================================

\begin{itemize}
    \item Cost estimates are based on publicly available API prices and may
          change; all raw token counts are logged to enable re-costing.
    \item HumanEval evaluation uses a non-execution-based heuristic in the
          current release; full functional correctness requires a sandboxed
          execution environment.
    \item The GAIA Level-1 subset excludes tasks that require tool use;
          Levels 2 and 3 are left to future work.
    \item All experiments use API-based inference; open-weight models are
          not yet supported.
\end{itemize}

% ============================================================
\section{Conclusion}
% ============================================================

We presented Pareto-Bench, a cost-controlled benchmark for comparing
multi-agent debate strategies against single-agent baselines.  By enforcing an
iso-cost protocol and drawing Pareto frontiers, we enable researchers to ask
whether debate architectures are genuinely worth their additional inference
budget---a question that existing benchmarks leave unanswered.  The full
codebase, data, and model registry are released at
\url{https://github.com/your-org/pareto-bench}.

\bibliographystyle{plainnat}
\bibliography{references}

\end{document}
